\BourbakiExerciseGroup

\begin{enumerate}
\def\labelenumi{\arabic{enumi})}
\item
  把命题 3、6、7、9 以及命题 4 的推论推广到半拓扑群（§ 1，习题 2）与仿拓扑群（§ 1，习题 4）。
\item
  \begin{enumerate}
  \def\labelenumii{\alph{enumii})}
  \tightlist
  \item
    把命题 1、2、4 推广到半拓扑群（§ 1，习题 2）。（为证明闭包 \(\overline{\mathbf{H}}\) （子群 \(\mathbf{H}\) 的）是子群，先注意 \(\overline{s\overline{\mathbf{H}}} \subset \overline{\mathbf{H}}\) 对一切 \(s \in \mathbf{H}\) 成立。）
  \end{enumerate}
\end{enumerate}

\begin{enumerate}
\def\labelenumi{\alph{enumi})}
\setcounter{enumi}{1}
\tightlist
\item
  设 \(\mathbf{G}\) 是群 \(\mathbf{Z}^{(\mathbb{N})}\) ，并设 \(e_i\) （ \(i \in \mathbb{N}\) ）是 \(\mathbf{G}\) 的这样的元素：除第 \(i\) 个坐标外所有坐标均为零，而第 \(i\) 个坐标为 1。对每个 \(n \in \mathbb{N}\) ，设 \(\mathbf{V}_n\) 是所有 \(z = (z_k) \in \mathbf{G}\) （满足 \(z_k = 0\) 对 \(k \leqslant n\) ，且 \(z_k \geqslant 0\) 对 \(k > n\) ）组成的集合。诸 \(\mathbf{V}_n\) 构成 \(0\) 的邻域基本系（对应于使 \(\mathbf{G}\) 成为豪斯多夫仿拓扑群的一个拓扑）。设 H 是 G 的由元素 \(e_{0} + e_{n} (n \geqslant 1)\) 生成的子群。试证 H 是离散的但在 G 中不闭，且 \(\overline{H}\) 不是 G 的子群。
\end{enumerate}

\begin{enumerate}
\def\labelenumi{\arabic{enumi})}
\setcounter{enumi}{2}
\tightlist
\item
  给出一个非豪斯多夫拓扑群的例子，其中中心不闭且中心的闭包是一个非交换子群（参看 § 1，习题 1）。给出一个具有同样性质的半拓扑群的例子，其中每个由单点组成的集合都是闭的{[}参看 § 1，习题 2 b){]}。
\end{enumerate}

¶ 4) a) 在半拓扑群 G 中，既开又闭的 e 的诸邻域的交是一个闭子群 H ，它在 G 的所有自同构下不变（注意不可能把 G 分划成两个既开又闭且都与 H 相交的子集）。

\begin{enumerate}
\def\labelenumi{\alph{enumi})}
\setcounter{enumi}{1}
\tightlist
\item
  在 G 的所有子群组成的集合 \(\mathfrak{G}\) 上，如下定义映射 \(\varphi: \mathfrak{G} \to \mathfrak{G}\) ：若 H 是 G 的任一子群，则 \(\varphi(H)\) 是 H 的子群，它是 H 中 \(e\) 的所有既开又闭的邻域的交。在集合 \(\mathfrak{G}\) （按关系 \(\supset\) 排序）中，考虑 G 的链 \(\Gamma\) （关于映射 \(\varphi\) ） （《集合论》，第三章，§ 2，习题 6）；试证 \(\Gamma\) 的最小子群是 \(e\) 在 G 中的连通分量（注意这个最小子群必定是连通的）。
\end{enumerate}

\begin{enumerate}
\def\labelenumi{\arabic{enumi})}
\setcounter{enumi}{4}
\tightlist
\item
  半拓扑群 G 是拟拓扑的，如果对每个 \(a \in G\) ，G 到 G 中的映射 \(x \to xax^{-1}\) 是连续的。每个交换半拓扑群都是拟拓扑的。
\end{enumerate}

\begin{enumerate}
\def\labelenumi{\alph{enumi})}
\item
  试证：拟拓扑群中闭子群的正规化子是闭的。
\item
  设 G 是一个拟拓扑群，它由 e 的每个邻域生成（相应地，它是连通的）。试证 G 的每个离散（相应地，全不连通）正规子群 D 都含于 G 的中心（若 \(a \in D\) ，试证存在 e 的一个邻域 V ，使得 \(xax^{-1} = a\) 对每个 \(x \in V\) 成立）。
\end{enumerate}

\(^{*}c)\) 设 \(G\) 是形如 \(\begin{pmatrix} 1 & 0 \\ y & x \end{pmatrix}\) 的矩阵所成的群，其中 \(x\) 与 \(y\) 是实数且 \(x > 0\) 。若把 \(G\) 与 \(R^2\) 的一个子集借助映射 \((x, y) \to \begin{pmatrix} 1 & 0 \\ y & x \end{pmatrix}\) 等同起来，则拓扑 \(\mathfrak{G}_0\) （在 \(G\) 上由 \(R^2\) 的拓扑诱导的）与 \(G\) 的群结构相容。设 \(H\) 是 \(G\) 的正规子群，由所有矩阵 \(\begin{pmatrix} 1 & 0 \\ y & 1 \end{pmatrix}\) 组成，并设 \(H^*\) 是 \(H\) 中单位元的余集。如下定义拓扑 \(\mathfrak{G}\) （比 \(\mathfrak{G}_0\) 更细）：取 \(e\) 在 \(G\) 中的邻域基本系为 \(e\) 的诸邻域（拓扑 \(\mathfrak{G}_{0}\) 中的）与 H* 在 G 中的余集之交。关于这个拓扑 \(\mathfrak{G}\) ，试证 G 是一个连通的半拓扑群，它不是拟拓扑的，且以 H 为一个不含于中心的离散子群。 *

* 6) 设 G 是正交群 \(\mathbf{O}_{2}(\mathbf{Q})\) ，把它与空间 \(\mathbf{Q}^{4}\) （即所有 \(2 \times 2\) 矩阵（系数在 \(\mathbf{Q}\) 中）组成的空间）的一个子空间等同起来。试证 \(\mathbf{Q}^{4}\) 的拓扑在 G 上诱导的拓扑与 G 的群结构相容。试证 G 的换位子群不闭，且它在 G 中的闭包是群 \(\mathbf{SO}_{2}(\mathbf{Q})\) 。 *

\begin{enumerate}
\def\labelenumi{\arabic{enumi})}
\setcounter{enumi}{6}
\tightlist
\item
  \begin{enumerate}
  \def\labelenumii{\alph{enumii})}
  \tightlist
  \item
    试证：若 G 是连通的拟拓扑群（习题 5），则 G 的换位子群是连通的。（设 \(P_k\) 是 k 个换位子 \(x^{-1}y^{-1}xy\) 的全部乘积组成的集合；试证 \(P_k\) 是一些连通集的并，其中每个都与 \(P_{k-1}\) 相交。）
  \end{enumerate}
\end{enumerate}

\begin{enumerate}
\def\labelenumi{\alph{enumi})}
\setcounter{enumi}{1}
\tightlist
\item
  设 G 是一个非交换无限群，其换位子群是有限的（例如，一个非交换有限群与一个无限交换群的积）。若 G 赋有 § 1 习题 2 b) 中定义的拓扑，试证 G 是一个连通的半拓扑群，其换位子群不连通。
\end{enumerate}

¶ 8) 设 G 是一个拟拓扑群（习题 5）。

\begin{enumerate}
\def\labelenumi{\alph{enumi})}
\item
  把命题 8 推广到 G （先证明 \(x\overline{\mathbf{K}}x^{-1} \subset \overline{\mathbf{K}}\) 对 \(x \in \mathbf{H}\) 成立）。给出一个半拓扑群的例子，它不是拟拓扑的且命题 8 对它不成立{[}参看 § 1，习题 2 b){]}。
\item
  设 H、K 是 G 的两个子群，使得 H⊃K 且 K 包含 H 的换位子群。试证 \(\overline{K}\) 包含 \(\overline{H}\) 的换位子群（类似的方法）。
\item
  由 b) 推断：若 G 的每个由单点组成的子集都是闭的，则 G 中每个可解（相应地，交换）子群的闭包是可解的（相应地，交换的）（就所考虑子群的导出列的长度进行论证）。
\end{enumerate}

\begin{enumerate}
\def\labelenumi{\arabic{enumi})}
\setcounter{enumi}{8}
\item
  设 G 是一个拟拓扑群，并设 H 是 G 的一个闭正规子群，它包含 G 的换位子群。试证：若 e 在 H 中的分量 K 是可解的，则 e 在 G 中的分量 L 也是可解的{[}利用习题 7 a) 证明 K 包含 L 的换位子群{]}。
\item
  设 G 是一个拟拓扑群，其中每个由单点组成的子集都是闭的，且其单位分量 C 使得 G/C 是有限的。试证：若 G/C 有 k 个元素，则共轭元 \(xax^{-1}\) （元素 \(a \in G\) 的）组成的集合或者是无限的，或者至多有 k 个元素（考虑 G 到这个集合上的映射 \(x \rightarrow xax^{-1}\) ）。特别地，若 G 是连通的，则不在 G 的中心的每个元素的共轭元集合是无限的。
\item
  设 G 是一个群（拓扑化与否均可），它作用于一个拓扑空间 X ，使得每个映射 \(x \rightarrow s.x\) （ \(s \in G\) ）（X 到 X 中的）都是连续的。把第 4 节的结果推广到这种情形。
\item
  设 G 是一个半拓扑（相应地，仿拓扑、拟拓扑）群，并设 H 是 G 的一个正规子群。
\end{enumerate}

\begin{enumerate}
\def\labelenumi{\alph{enumi})}
\item
  试证：若 G/H 赋以 G 关于 H 的商拓扑，则 G/H 是一个半拓扑（相应地，仿拓扑、拟拓扑）群。
\item
  把第 7 节的结果推广到半拓扑（相应地，仿拓扑、拟拓扑）群。
\end{enumerate}

\begin{enumerate}
\def\labelenumi{\arabic{enumi})}
\setcounter{enumi}{12}
\item
  设 G 是一个半拓扑（相应地，仿拓扑）群，并设 H 是 G 的一个稠密子群。齐性空间 G/H 的拓扑是什么？
\item
  设 G 是一个半拓扑群，并设 H 是 G 的一个正规子群。试证半拓扑群 \(G/\overline{H}\) 同构于与 G/H 相伴的豪斯多夫群。
\end{enumerate}

¶ 15) 设 G 是一个拟拓扑群，其中每个由单点组成的集合都是闭的。试证：若 G 是可解的，则存在 G 的一个由闭子群组成的合成列，使得该列的逐次商群都是交换的{[}利用习题 8 c)，就 G 的导出列的长度用归纳法论证{]}。

¶ 16) 设 H 是半拓扑群 G 的一个子群，含于 e 的分量 K 中。试证：空间 G/H 的诸分量是 G 的诸分量在 G 到 G/H 上的典范映射 f 下的像{[}若 L 是 G/H 的一个分量，试证 \(\overline{f}^{-1}(L)\) 是连通的（用反证法）{]}。试证 K 是 G 的闭子群 H （使 G/H 全不连通者）中最小的一个。

* 17) 设 G 是 N 到 Q 中的映射 \(n \to f(n)\) （使得 \(\lim_{n \to \infty} f(n)\) 在 R 中存在）组成的加法群。对每个 \(\alpha > 0\) ，设 \(V_{\alpha}\) 是 \(f \in G\) （满足 \(|f(n)| < \alpha\) ，对一切 \(n \in \mathbb{N}\) ）组成的集合。试证诸 \(V_{\alpha}\) 满足公理 \((GV_{I})\) 与 \((GV_{II})\) ，且带由诸 \(V_{\alpha}\) 所定义的拓扑，群 G 是豪斯多夫的且全不连通。设 H 是所有 \(f \in G\) （满足 \(\lim_{n \to \infty} f(n) = 0\) ）组成的子群。试证 H 是闭的，且 G/H 同构于 R ，因而是连通的。 *

\begin{enumerate}
\def\labelenumi{\arabic{enumi})}
\setcounter{enumi}{17}
\item
  一个连续同态 \(f\) （拓扑群 \(G\) 到拓扑群 \(G'\) 中的）是严格态射，当且仅当每个开集在 \(f\) 下的像（开集取自 \(G\) ）相对于 \(f(G)\) 至少有一个内点。
\item
  设 \(G, G', G''\) 是三个拓扑群；设 \(f\) 是 \(G\) 到 \(G'\) 中的一个严格态射，并设 \(g\) 是 \(F'\) 到 \(G''\) 中的一个严格态射。试证：若 \(f(G)\) 包含核 \(\overline{g}^{-1}(e'')\) （\(g\) 的，\(e''\) 是 \(G''\) 的单位元），则 \(g \circ f\) 是 \(G\) 到 \(G''\) 中的一个严格态射（参看第 7 节，命题 21 的推论）。给出一个例子，其中 \(f\) 是单射的且前一条件不满足，而 \(g \circ f\) 不是 \(G\) 到 \(G''\) 中的严格态射（参看第 7 节，命题 20 后面的注）。
\item
  设 H、H' 是半拓扑群 G 的两个子群，使得 \(H' \subset H\) 。试证：若 \(f, f'\) 分别是 G 到 G/H 与 \(G/H'\) 上的典范映射，则存在唯一的映射 \(\varphi\) （\(G/H'\) 到 G/H 上的），使得 \(f = \varphi \circ f'\) ；\(\varphi\) 是连续的且是开的。此外若 \(H'\) 在 H 中是开的，则 \(G/H'\) 的每一点 x 有一个邻域 V ，使得 \(\varphi\) 在 V 上的限制是 V 到 G/H 中 \(\varphi(x)\) 的一个邻域上的同胚。
\item
  设 G 是一个拓扑群，并设 H、K 是 G 的两个子群，使得 G = HK 且 \(H \cap K = \{e\}\) 。每个 \(x \in G\) 都可以唯一地写成 \(x = f(x)g(x)\) 的形式，其中 \(f(x) \in H\) 且 \(g(x) \in K\) 。映射 \((y, z) \to yz\) （\(H \times K\) 到 G 上的）是同胚，当且仅当映射 f、g 之一是连续的。试证：若子群 H、K 之一是紧的且闭的，而另一个是闭的，则这个条件总满足（注意若 H 与 K 都是闭的，则当 x 在 G 中趋于 e 时，f 与 g 除 e 外不能有任何聚点）。
\item
  设 G 是拓扑群的一个族 \((\mathbf{G}_{i})_{i\in I}\) 的积，并设 \(G'\) 是一个拓扑群，使得存在一个邻域 \(V'\) （单位元 \(e'\) 在 \(G'\) 中的），它不含 \(G'\) 的除 \(\{e'\}\) 外的任何正规子群。若 f 是 G 到 \(G'\) 中的一个连续同态，试证存在 I 的一个子集 J ，其余集 \(J'\) 是有限的，使得 f 等于 \(e'\) （在 G 的子群 \(\prod_{i\in I}G_{i}\times\prod_{i\in J'}\{e_{i}\}\) 上）。
\end{enumerate}

\[
\left(\mathbf {G} _ {i}\right) _ {i \in \mathbf {I}}
\]

\[
\mathbf {G} _ {i}
\]

\[
\prod_ {i \in I} A _ {i},
\]

其中 \(\mathbf{A}_i\) 在 \(\mathbf{G}_i\) 中开（对每个 \(i \in I\)），且 \(i \in I\) （使 \(\mathbf{A}_i \neq \mathbf{G}_i\) 者）组成的集合具有基数 \(< c\) ，这里 \(c\) 是一个给定的无限基数（第一章，§ 4，习题 9）。试证这个拓扑与 G 的群结构相容。由此构造一个非离散豪斯多夫拓扑群的例子，其中每个紧子集都是有限的（参看第一章，§ 9，习题 4）。

\begin{enumerate}
\def\labelenumi{\arabic{enumi})}
\setcounter{enumi}{23}
\tightlist
\item
  设 G 与 G' 是两个局部同构的连通群。
\end{enumerate}

\begin{enumerate}
\def\labelenumi{\alph{enumi})}
\item
  设 \(f\) 是 \(\mathbf{G}\) 与 \(\mathbf{G}'\) 之间的一个局部同构，定义在 \(\mathbf{V}\) （\(\mathbf{G}\) 的单位元的一个邻域）上。设 \(\mathbf{H}\) 是积 \(\mathbf{G} \times \mathbf{G}'\) 的由点 \((x, f(x))\) （\(x \in \mathbf{V}\)）所成的集生成的子群。考虑 \(\mathbf{H}\) 上的滤基，它由 \(e\) 的邻域（含于 \(\mathbf{V}\) 者）在映射 \(x \to (x, f(x))\) 下的像组成。试证这个滤基是单位元（\(\mathbf{H}\) 的）的邻域基本系（对应于一个拓扑 \(\mathcal{C}\) ，它与 \(\mathbf{H}\) 的群结构相容）。
\item
  试证存在两个离散正规子群 K、K' 含于 H 的中心，使得 G 同构于 H/K ，且 G' 同构于 H/K' {[}利用习题 5 b){]}。
\end{enumerate}

* c) 取 G 与 G' 都为群 T ；设 φ 是 R 到 T = R/Z 上的典范同态，并设 f 是这样一个局部同构：它把每个点 φ(x) （其中 \(-1/4 \leqslant x \leqslant 1/4\) ）映为 φ(θx) ，θ 是一个无理数且 0 \textless{} θ \textless{} 1。试证对这个局部同构，在 H 上定义的拓扑 C 不同于由 G × G' 上积拓扑在 H 上诱导的拓扑*。

\begin{enumerate}
\def\labelenumi{\arabic{enumi})}
\setcounter{enumi}{24}
\item
  设 G 是一个拓扑群，K 是 G 的一个正规子群，\(\varphi\) 是典范映射 \(G \rightarrow G/K\) 。设 f 是拓扑群 H 到 G 中的一个连续同态，使得 \(f(H)\) 的每个元素与 K 的每个元素交换。则 \((y, z) \rightarrow f(y)z\) 是 \(H \times K\) 到 G 中的一个连续同态 g 。g 是 \(H \times K\) 到 G 中的一个严格态射，当且仅当 \(\varphi \circ f\) 是 H 到 G/K 中的一个严格态射。
\item
  设 \((\mathbf{G}_i)_{i\in I}\) 是拓扑群的一个族，且对每个 \(i\in I\) ，设 \(\mathbf{K}_i\) 是 \(\mathbf{G}_i\) 的一个开正规子群。用 \(\mathfrak{B}\) 表示积拓扑群 \(\mathbf{K} = \prod_{i\in I}\mathbf{K}_i\) 中单位元的邻域滤子。试证：在积群 \(\mathbf{G} = \prod_{i\in I}\mathbf{G}_i\) 中，\(\mathfrak{B}\) 是 \(e\) 的邻域滤子的一个基（对应于一个与 \(\mathbf{G}\) 的群结构相容的拓扑）。群 \(\mathbf{G}\) （带此拓扑）称为诸 \(\mathbf{G}_i\) 的局部积（相对于诸 \(\mathbf{K}_i\) ）；此时 \(\mathbf{K}\) 是 \(\mathbf{G}\) 的一个开正规子群，且 \(\mathbf{G} / \mathbf{K}\) 同构于诸群 \(\mathbf{G}_i / \mathbf{K}_i\) 的（离散）积。若对每个 \(i\in I\) ，\(\mathbf{K}_i'\) 是 \(\mathbf{G}_i\) 的另一个开正规子群，则 \(\mathbf{G}\) 上由族 \((\mathbf{K}_i)\) 与 \((\mathbf{K}_i')\) 定义的局部积拓扑彼此一致，当且仅当 \(\mathbf{K}_i' = \mathbf{K}_i\) 对除有限多个指标外的所有指标成立。
\end{enumerate}

对每个 \(i \in I\) ，设 \(G_i'\) 是 \(G\) 的正规子群，由所有 \(x = (x_i)\) （满足 \(x_i = e_i\) ，只要 \(i \neq x\) ）组成。则 \(G_i'\) （带由 \(G\) 上一个局部积拓扑诱导的拓扑）典范地同构于 \(G_i\) 。在 \(G\) 中，由诸 \(G_i'\) 生成的子群的闭包是正规子群 \(G_0\) （\(G\) 的），它由所有 \(x = (x_i)\) （满足 \(x_i \in K_i\) ，对除有限多个指标外的所有指标）组成。\(G_0\) 称为诸 \(G_i\) 的局部直积（相对于诸 \(K_i\) ）；此时 \(G_0 / K\) 是一个离散群，同构于诸 \(G_i / K_i\) 的直积。

¶ 27) 称群 G 具有性质 P(n) （n 是一个整数 \textgreater0），如果对 G 的 n 个元素的每个组 \((s_{1}, \ldots, s_{n})\) ，存在一个置换 \(\sigma \in S_{n}\) ，它异于恒等置换（但依赖于 \(s_{1}, \ldots, s_{n}\) ），使得 \(s_{\sigma(1)} \ldots s_{\sigma(n)} = s_{1} \ldots s_{n}\) 。

\begin{enumerate}
\def\labelenumi{\alph{enumi})}
\item
  试证非交换群 \(\mathfrak{S}_3\) 满足 \(\mathbf{P}(4)\) 。
\item
  设 G 是一个连通豪斯多夫拓扑群。试证：若 G 满足 \(\mathrm{P}(n)\) （对某个整数 n \textgreater{} 1），则 G 是交换的。{[}试证 G 满足 \(\mathrm{P}(n - 1)\) ：给定 G 的 n - 1 个元素 \(s_{1}, \ldots, s_{n-1}\) ，归结到存在 e 的一个邻域 U 的情形，使得对每个 \(x \in U\) ，有 \(s_{1} \ldots s_{n-1}x = s_{1} \ldots s_{i}xs_{i+1} \ldots s_{n-1}\) （对某个指标 \(i = i(x)\) ）；推出元素 \(s_{j+1} \ldots s_{n-1} (0 \leqslant j \leqslant n - 2)\) 之一的中心化子在 G 中既开又闭，从而等于 G。于是当 \(j \geqslant 1\) 时结论得证；若 j = 0，试证
\end{enumerate}

\[
\left. s _ {1} s _ {2} \dots s _ {n - 1} = s _ {2} \dots s _ {n - 1} s _ {1} \right].
\]

¶ 28) a) 设 S 是拓扑群 G 的一个稳定子集。试证 \(\mathring{S}\) 与 \(\overline{S}\) 都是稳定的，且 \(SS\subset\mathring{S}\) 与 \(\mathring{SS}\subset\mathring{S}\) 。若此外 S 是开的且 \(e\in\overline{S}\) ，则 \(S=\mathring{\overline{S}}\) （若 \(x\in\overline{S}-S\) ，试证 \(x\) 是 S 的一个边缘点）。

\begin{enumerate}
\def\labelenumi{\alph{enumi})}
\setcounter{enumi}{1}
\item
  在本题的其余部分，G 是一个交换拓扑群，写成加法。对 G 的每个非空子集 A ，用 \(s(A)\) 表示所有 \(x \in G\) （满足 \(x + A \subset A\) 者）组成的集合，即诸集合 A - y （y 取遍 A）的交；设 \(b(A) = s(A) \cap s(-A)\) 。集合 \(s(A)\) 是 G 的一个稳定子集，而 \(b(A)\) 是 G 的一个子群，由所有 \(x \in G\) （满足 \(x + A = A\) 者）组成。试证：若 F 是 G 的任一非空闭子集，则 \(s(F)\) 是闭的。对 G 的每个非空子集 A ，我们有 \(s(A) \subset s(\overline{A})\) 与 \(s(A) \subset s(\mathring{A})\) ；而若 \(A = \mathring{\overline{A}}\) ，则 \(s(\mathring{\overline{A}}) = s(\overline{\overline{A}})\) 。
\item
  设 \(\mathfrak{S}\) 是所有稳定开子集 \(S\) （\(G\) 的，满足 \(0 \notin S\) 者）组成的集合，并设 \(\mathfrak{S} \neq \emptyset\) 。则集合 \(\mathfrak{S}\) （按包含关系排序）是归纳的；设 \(\mathfrak{M} \neq \emptyset\) 是其极大元组成的集合。试证：若 \(\mathbf{M} \in \mathfrak{M}\) ，则 \(\mathbf{M} = \overline{\mathbf{M}}\) {[}利用 \(a\) {]}，且对每个整数 \(n > 0\) ，\(\mathbf{M}\) 等于所有 \(x \in \mathbf{G}\) （满足 \(nx \in \mathbf{M}\) 者）组成的集合。此外，子群 \(b(\mathbf{M})\) 等于 \(\mathbf{M} \cup (-\mathbf{M})\) 在 \(\mathbf{G}\) 中的余集（若 \(x \notin \mathbf{M}\) 且 \(-x \notin \mathbf{M}\) ，考虑诸集合 \(kx + \mathbf{M}\) （整数 \(k \geqslant 0\) ）的并）；推出 \(\mathbf{G} = \mathbf{M} - \mathbf{M}\) 。
\end{enumerate}

一个闭子群 \(\mathbf{B}\) （\(\mathbf{G}\) 的）称为剩余的，如果存在 \(\mathbf{M} \in \mathfrak{M}\) 使得 \(\mathbf{B} = b(\mathbf{M})\) 。

\begin{enumerate}
\def\labelenumi{\alph{enumi})}
\setcounter{enumi}{3}
\item
  G 的所有剩余子群的交称为 G 的根，记作 \(T_{G}\) 。若 \(T_{G}=G\) （相应地，\(T_{G}=\{0\}\) ），则称 G 是一个根群（相应地，无根群）。x 属于 \(T_{G}\) ，当且仅当 G 的每个包含 x 的稳定开子集也包含 0。由此推出：若 G 是离散的，则 \(T_{G}\) 是 G 的挠子群。若 f 是 G 到交换拓扑群 \(G'\) 中的任一连续同态，则 \(f(T_{G})\subset T_{G'}\) 。若 H 是 G 的一个含于 \(T_{G}\) 的子群，则 \(T_{G/H}=T_{G}/H\) 。子群 \(T_{G}\) 是 G 的子群 H （使 G/H 无根者）中最小的一个。
\item
  试证：在 G 的满足 \(T_{H} = H\) 的诸子群 H 中，有一个最大的 \(H_{0} \subset T_{G}\) ，它是闭的。
\end{enumerate}

\(f)\) 若 \(\mathbf{T}_{\mathbf{G}} \neq \mathbf{G}\) ，则 \(\mathbf{T}_{\mathbf{G}}\) 是开的，当且仅当 \(\mathbf{G}\) 中存在一个剩余的开子群{[}试证：若 \(b(\mathbf{M})\) 对某个 \(\mathbf{M} \in \mathfrak{M}\) 是开的，则 \(\mathbf{T}_{\mathbf{G}} \supset b(\mathbf{M})\) （利用 \(e)\) ）{]}。

* 29) 设 G 是拓扑群 R ，\(\varphi\) 是 R 到其商群 \(\mathbf{T} = \mathbf{R} / \mathbf{Z}\) 上的典范映射，\(\theta\) 是一个无理数。群 G 连续地作用于 \(\mathbf{E} = \mathbf{T}^2\) ，按规则 \((s, (x, y)) \to (x + \varphi(s), y + \varphi(\theta s))\) 。试证每个 \(z \in \mathbf{E}\) 的稳定子群仅由 0 组成，\(z\) 的轨道在 E 中稠密，且它不是相对于 G 的一个拓扑齐性空间*。

\begin{enumerate}
\def\labelenumi{\arabic{enumi})}
\setcounter{enumi}{29}
\tightlist
\item
  称拓扑群 G 没有任意小的子群，如果存在 e 在 G 中的一个邻域 V ，使得 \(\{e\}\) 是 G 的含于 V 的唯一子群。
\end{enumerate}

\begin{enumerate}
\def\labelenumi{\alph{enumi})}
\item
  设 G 是一个拓扑群，H 是 G 的一个正规子群。试证：若 H 与 G/H 都没有任意小的子群，则 G 也没有。
\item
  由 a) 推断：若 \(H_{1}\) 、\(H_{2}\) 是拓扑群 G 的两个正规子群，使得 \(G/H_{1}\) 与 \(G/H_{2}\) 都没有任意小的子群，则 \(\mathrm{G}/(\mathrm{H}_{1} \cap \mathrm{H}_{2})\) 也没有。
\end{enumerate}

¶ 31) 设 G 是一个连通拓扑群，H 是 G 的一个子群，U 是 G 的一个开子集，使得 H = UH。设 \(V = H \cap (U^{-1}U)\) ；试证 \(H = V^{\infty}\) {[}试证：若 \(A = H \cap C(V^{\infty})\)，则集合 \(U.V^{\infty}\) 与 \(U.A\) 不相交{]}。

